\chapter{Introduction}

\section{General}
Agricultural markets (\textit{mandis}) in India are numerous, fragmented,
and geographically dispersed. A single commodity can trade at meaningfully
different prices at mandis only a few hours apart on the same day, yet the
individual farmer -- especially small and marginal farmers who form the
majority of India's agricultural workforce -- typically has no easy way to
compare these prices before deciding where and when to sell. This
information asymmetry, combined with limited on-farm storage capacity,
frequently forces farmers into \textit{distress selling}: selling
immediately after harvest, at the point of lowest leverage, often to
intermediaries who pay below the prevailing market rate.

This project, the \textbf{\ProjectTitle}, addresses this gap directly. It
combines publicly available government price data (Agmarknet, under the
Ministry of Agriculture \& Farmers Welfare, Government of India) with a
time-series forecasting pipeline and a farmer-facing web application to
give users a live, comparative, and forward-looking view of mandi prices
for the crops and markets relevant to them.

\section{Problem Statement}
Given the absence of an accessible, forecast-aware, mandi-comparison tool
for Indian farmers, this project aims to answer the following question:
\textit{Can publicly available agricultural market data be used to build a
system that (a) forecasts short-term commodity prices at the level of
individual mandis, and (b) presents this information, along with an
actionable Sell/Hold recommendation, through an interface usable by a
non-technical farmer?}

\section{Objectives}
The specific objectives of this project are:
\begin{enumerate}[label=\arabic*.]
    \item To identify and consolidate reliable historical and live sources
    of Indian agricultural commodity price data (Agmarknet, CEDA Agri
    Market Data).
    \item To design a data pipeline that ingests live daily price data on
    an ongoing basis while maintaining sufficient historical depth for
    time-series modelling.
    \item To build a short-term price forecasting model at the granularity
    of individual (commodity, mandi) pairs, rather than a blurred
    state- or commodity-level average, so that local price variation --
    the core value proposition of the system -- is preserved.
    \item To translate model output into a simple, actionable Sell/Hold
    signal that a farmer can interpret without statistical background.
    \item To design and implement a two-page web application: a
    general dashboard (state-wise price overview) and an interactive
    prediction page allowing state- and crop-level drill-down to
    individual mandi comparisons.
    \item To evaluate the forecasting approach against a naive baseline
    and document the comparison for academic assessment.
\end{enumerate}

\section{Scope of the Project}
The system, in its current planned scope, covers 2--3 Indian states and
the top six most-consistently-reported commodities per state (selected
using actual historical reporting-frequency data rather than assumption),
across roughly 15--20 mandis per state. This scope is intentionally
bounded for a single academic-year capstone project; the architecture is
designed to extend to additional states/commodities without structural
change (see Chapter~3).

\section{Organization of the Report}
Chapter~2 reviews existing approaches to agricultural price prediction and
related student/academic projects, and identifies the gap this project
addresses. Chapter~3 presents the proposed system architecture, data
pipeline, machine learning design, and technology stack, along with a
feasibility discussion. Chapter~4 will present implementation results and
evaluation (populated in subsequent review periods as development
progresses). Chapter~5 presents the current work plan, timeline, and
identifies the project's present limitations and future scope.
